% ---------------------------------------------------------------------------
% Chapter: Toolkit overview and parakit
% ---------------------------------------------------------------------------
\chapter{Toolkit Overview and Parameter Management}\label{chap:toolkit}

The \pkg{toolkit} package collects research utilities that are independent of
the environments. This chapter describes its structure and then documents
\pkg{toolkit.parakit}, the parameter-management module. The plotting and
neural-network subpackages have their own chapters.

\section{Structure of the toolkit}\label{sec:toolkit-structure}

\begin{table}[htbp]
  \centering
  \small
  \caption{Subpackages of \pkg{toolkit}.}
  \label{tab:toolkit-subpackages}
  \begin{tabularx}{\textwidth}{@{}l>{\raggedright\arraybackslash}X>{\raggedright\arraybackslash}p{3.3cm}l@{}}
    \toprule
    Subpackage & Contents & Dependencies & Chapter \\
    \midrule
    \pkg{toolkit.plotkit} & Learning curves, bars, histograms, heatmaps,
      style presets, palettes, figure export, gallery CLI &
      matplotlib (core) & \ref{chap:plotkit} \\
    \pkg{toolkit.neural\_toolkit} & Policy, value and Q-networks, encoders,
      decoders, factories, \code{NetworkUtils}, tabular tools &
      PyTorch (\code{torch} extra) & \ref{chap:neural} \\
    \pkg{toolkit.parakit} & Save, load, validate and apply \pkg{argparse}
      defaults; Tk editor window & standard library; Tk for the window &
      \ref{chap:toolkit} \\
    \bottomrule
  \end{tabularx}
\end{table}

The subpackages are imported lazily: \code{import toolkit} only defines
\code{toolkit.\_\_version\_\_}, and the first access to
\code{toolkit.plotkit}, \code{toolkit.neural\_toolkit} or
\code{toolkit.parakit} imports that subpackage. A script that only plots
therefore never imports PyTorch, and \pkg{parakit}'s headless functions work
on systems without Tk. The usual import style names the subpackage
explicitly:

\begin{pycode}
from toolkit.plotkit import plot_learning_curves, save_figure
from toolkit.parakit import apply_parameters, load_parameters
from toolkit.neural_toolkit import MLPPolicyNetwork, NetworkUtils  # needs torch
\end{pycode}

All three subpackages follow the conventions of the environments: they never
print, never configure logging (diagnostics go to \code{logging} loggers named
after the module), never touch the global random state, and report problems
with exceptions or \code{warnings} whose messages say how to fix them.

\section{parakit: managing argparse parameters}\label{sec:parakit}

Research scripts are usually configured through an
\code{argparse.ArgumentParser}. \pkg{parakit} treats the parser as the single
source of truth for the parameters, their types and their allowed values, and
adds four operations on top of it: reading the current defaults, saving them
to a JSON file, loading a file, and installing validated values as the new
defaults. Because the values become \emph{defaults}, flags given on the
command line still override them. An optional Tk window
(Section~\ref{sec:parakit-gui}) lets users edit the values interactively.

\subsection{Quick start}\label{sec:parakit-quickstart}

The typical workflow saves the configuration of a run next to its results and
later restores it:

\begin{pycode}
import argparse
from toolkit.parakit import (
    apply_parameters, get_parameters, load_parameters, save_parameters,
)

parser = argparse.ArgumentParser()
parser.add_argument("--lr", type=float, default=3e-4, help="learning rate")
parser.add_argument("--batch_size", type=int, default=64)
parser.add_argument("--optimizer", choices=["adam", "sgd"], default="adam")
parser.add_argument("--hidden", type=int, nargs="+", default=[256, 256])
parser.add_argument("--layer_norm", action="store_true")
parser.add_argument("--checkpoint", default=None)

# Save the current defaults into a directory (timestamped file name).
path = save_parameters(get_parameters(parser), "runs/params")

# Later: load the newest file of the directory and install it as defaults.
apply_parameters(parser, load_parameters("runs/params"))

# Values given as text are converted according to the parser.
apply_parameters(parser, {"lr": "1e-3", "hidden": "64 64", "layer_norm": "yes"})
args = parser.parse_args([])            # lr=0.001, hidden=[64, 64], layer_norm=True
args = parser.parse_args(["--lr", "0.01"])  # command-line flags still win
\end{pycode}

\subsection{Headless API}\label{sec:parakit-api}

Table~\ref{tab:parakit-api} lists the functions; none of them needs Tk.

\begin{table}[htbp]
  \centering
  \small
  \caption{Headless functions of \pkg{toolkit.parakit}.}
  \label{tab:parakit-api}
  \begin{tabularx}{\textwidth}{@{}>{\raggedright\arraybackslash}p{5.6cm}>{\raggedright\arraybackslash}X@{}}
    \toprule
    Function & Description \\
    \midrule
    \code{get\_parameters(parser)} &
      Dictionary \code{\{dest: default\}} of all tunable options (lists are
      copied). \\
    \code{save\_parameters(values, path\_or\_dir)} &
      Write \code{values} as JSON and return the \code{Path} of the file.
      A path with a suffix is used as the file name; a directory (existing, or
      a path without suffix) receives a new timestamped file. Parent
      directories are created and the write is atomic. \\
    \code{load\_parameters(path)} &
      Read a JSON object from a file, or from the newest
      \code{parameters\_*.json} file of a directory. Raises
      \code{FileNotFoundError} or \code{ValueError}. \\
    \code{apply\_parameters(parser, values, strict=True, validation\_callbacks=None)} &
      Convert and validate every value with \code{parse\_value}, then call
      \code{parser.set\_defaults}. Returns the parser. \\
    \code{parse\_value(action, text, validator=None)} &
      Convert one value (text or an already typed JSON value) for an
      \pkg{argparse} action; see Section~\ref{sec:parakit-parsing}. \\
    \code{parse\_bool(value)} &
      Convert \code{true/false}, \code{yes/no}, \code{on/off}, \code{1/0}
      (any case) or a boolean to \code{bool}. \\
    \code{format\_value(action, value)} &
      Text that \code{parse\_value} converts back (lists as JSON, \code{None}
      as \code{"None"}). \\
    \code{describe\_type(action)} &
      Short type description such as \code{"float"}, \code{"list[int]"} or
      \code{"str or None"}. \\
    \code{tunable\_actions(parser)} &
      Dictionary \code{\{dest: action\}} of the options \pkg{parakit}
      manages. \\
    \bottomrule
  \end{tabularx}
\end{table}

\paragraph{Tunable options.}
\code{tunable\_actions} skips help, version and sub-parser actions and options
whose destination or default is \code{argparse.SUPPRESS}. When several options
share a destination, such as a \code{--flag}/\code{--no-flag} pair, the first
one is used. Positional arguments are included and keyed by their name.

\paragraph{Errors.}
\code{apply\_parameters} either applies all values or none. With
\code{strict=True} (the default) a key that is not a parser destination raises
a \code{ValueError} that lists the unknown keys and the valid names; with
\code{strict=False} unknown keys are ignored. Otherwise every value is
converted, and all conversion and validation failures are reported together
in one \code{ValueError}, one line per parameter. The parser is left unchanged
whenever an exception is raised.

\subsection{File format}\label{sec:parakit-format}

A parameter file is a flat JSON object that maps each destination name
(\code{action.dest}) to its value, written in UTF-8 with an indentation of four
spaces:

\begin{textcode}
{
    "lr": 0.001,
    "batch_size": 64,
    "optimizer": "adam",
    "hidden": [
        64,
        64
    ],
    "layer_norm": true,
    "checkpoint": null
}
\end{textcode}

\begin{itemize}
  \item \emph{File names.} Saving into a directory creates
    \code{parameters\_YYYYmmdd\_HHMMSS.json}; if that name exists (two saves in
    the same second), \code{\_1}, \code{\_2}, \ldots\ is appended. Loading from
    a directory reads the file whose name sorts last, which is the most recent
    one.
  \item \emph{Values.} JSON numbers, strings, booleans, \code{null} and lists
    are stored as such. NumPy arrays and scalars are converted with
    \code{tolist()}, sets are stored as sorted lists and any other object as its
    \code{str()}.
  \item \emph{Loading.} Values read from a file pass through the same
    conversion as text, so a file may contain \code{"64"} or \code{64} for an
    integer option, and \code{"yes"} or \code{true} for a flag. A file may
    contain a subset of the parameters; the other defaults are kept.
  \item \emph{Atomic writes.} Files are written to a temporary file in the
    target directory and renamed, so an interrupted save never leaves a
    truncated file behind.
\end{itemize}

\subsection{Parsing rules}\label{sec:parakit-parsing}

\code{parse\_value(action, text)} derives the conversion from the parser's
own declarations, so the same rules apply to values typed in the editor
window, loaded from JSON or passed to \code{apply\_parameters}.
Table~\ref{tab:parakit-rules} summarises them.

\begin{table}[htbp]
  \centering
  \small
  \caption{Conversion rules of \code{parse\_value}.}
  \label{tab:parakit-rules}
  \begin{tabularx}{\textwidth}{@{}>{\raggedright\arraybackslash}p{4.4cm}>{\raggedright\arraybackslash}X@{}}
    \toprule
    Declaration & Accepted input and result \\
    \midrule
    \code{type=callable} &
      The callable is applied to text (surrounding blanks are removed for
      \code{int} and \code{float}). Already typed values from JSON are checked:
      an \code{int} option accepts \code{3} and \code{3.0} but not \code{3.5};
      a \code{float} option accepts integers. \\
    No \code{type} &
      The type is inferred from the default (\code{bool}, \code{int},
      \code{float}, otherwise \code{str}), or from the \code{choices} when all
      have the same type. \\
    Booleans: \code{store\_true}, \code{store\_false},
      \code{BooleanOptionalAction}, \code{type=bool} or a boolean default &
      \code{true/false}, \code{yes/no}, \code{on/off}, \code{1/0} in any case,
      or a JSON boolean. The string \code{"False"} gives \code{False}. \\
    Lists: \code{nargs="+"}, \code{"*"}, an integer, \code{append} and
      \code{extend} actions &
      A JSON list (\code{[1.0, 2.0]}), a Python literal list or tuple, or
      tokens separated by commas (if the text contains a comma) or by
      whitespace. Each element is converted; \code{"+"} requires at least one
      element and an integer \code{nargs} exactly that many. \\
    Default \code{None} or \code{nargs="?"} &
      \code{None}, \code{null} (any case) and the empty string give
      \code{None}. For other options \code{None} is rejected, except that a
      string option accepts the text \code{"None"} literally. \\
    \code{choices} &
      Enforced on the converted value (on every element of a list). \\
    \code{count} action & Converted as an integer. \\
    \bottomrule
  \end{tabularx}
\end{table}

\paragraph{Declare types explicitly.}
Type inference from the default is a convenience of \pkg{parakit}, not of
\pkg{argparse}. An option declared as \code{add\_argument("--gamma",
default=0.99)} is converted to \code{float} by \pkg{parakit}, but
\code{parser.parse\_args(["--gamma", "0.9"])} returns the \emph{string}
\code{"0.9"}. Give every option a \code{type=} so that both paths agree.

\paragraph{Validation callbacks.}
A validator receives the \emph{converted} value and signals the result in one
of the following ways: return \code{True}/\code{False}; return a tuple
\code{(is\_valid, message)} or the legacy form \code{(is\_valid,
converted\_value, message)}; return \code{None} (valid); or raise
\code{ValueError} or \code{TypeError}, whose text becomes the error message.
Validators are not called when the converted value is \code{None}, and a value
returned by a validator is ignored: the result always follows the parser
declaration.

\begin{pycode}
def positive_even(value):
    return value > 0 and value % 2 == 0, "batch size must be a positive even number"

validators = {"batch_size": positive_even,
              "lr": lambda v: (0 < v <= 1, "lr must lie in (0, 1]")}
apply_parameters(parser, {"batch_size": "128", "lr": 0.01},
                 validation_callbacks=validators)

try:
    apply_parameters(parser, {"batch_size": 7, "optimizer": "rmsprop"},
                     validation_callbacks=validators)
except ValueError as exc:
    print(exc)
# Invalid parameter values:
# batch_size: batch size must be a positive even number
# optimizer: 'rmsprop' is not a valid choice; value must be one of: 'adam', 'sgd'
\end{pycode}

\subsection{Recommended workflow in training scripts}\label{sec:parakit-workflow}

The following pattern gives a script three configuration layers with a clear
precedence: parser defaults, then an optional parameter file, then explicit
command-line flags.

\begin{pycode}
import argparse
import sys
from toolkit.parakit import apply_parameters, get_parameters, load_parameters, save_parameters

def build_parser():
    parser = argparse.ArgumentParser()
    parser.add_argument("--config", default=None, help="JSON file or directory")
    parser.add_argument("--lr", type=float, default=3e-4)
    parser.add_argument("--seed", type=int, default=0)
    return parser

def main(argv=None):
    parser = build_parser()
    pre, _ = parser.parse_known_args(argv)          # read --config first
    if pre.config is not None:
        apply_parameters(parser, load_parameters(pre.config), strict=False)
    args = parser.parse_args(argv)                  # flags override the file
    run_dir = f"runs/seed{args.seed}"
    save_parameters(vars(args), f"{run_dir}/config.json")
    return args

main(["--lr", "0.001"])
\end{pycode}

\code{AJLATTConfig.add\_arguments(parser)} registers every AJLATT option on a
parser (Chapter~\ref{chap:ajlatt}), so environment parameters can be managed
in the same way.

\subsection{The ParameterTuner window}\label{sec:parakit-gui}

\code{ParameterTuner} opens a Tk window with one input field per tunable
option: a check box for booleans, a drop-down list for options with
\code{choices}, and a text field otherwise. Each field shows the option's help
text and its type as reported by \code{describe\_type}. Values are validated
with \code{parse\_value} (and the optional validation callbacks), written to
timestamped JSON files, and installed as the parser's defaults when the window
closes.

\begin{table}[htbp]
  \centering
  \small
  \caption{Constructor arguments of \code{ParameterTuner}.}
  \label{tab:parakit-tuner}
  \begin{tabularx}{\textwidth}{@{}>{\raggedright\arraybackslash}p{3.8cm}>{\raggedright\arraybackslash}p{2.1cm}>{\raggedright\arraybackslash}X@{}}
    \toprule
    Argument & Default & Meaning \\
    \midrule
    \code{parser} & required & The \code{argparse.ArgumentParser} to tune. \\
    \code{save\_delay} & \code{30000} & Auto-save interval in milliseconds. \\
    \code{save\_path} & \code{None} & Output directory for the JSON files; \code{None}
      uses the current working directory. Created on the first save. \\
    \code{auto\_save\_enabled} & \code{True} & Start with auto-save switched on. \\
    \code{validation\_callbacks} & \code{None} & Per-destination validators
      (Section~\ref{sec:parakit-parsing}); callbacks for unknown names are
      ignored with a warning. \\
    \code{inactivity\_timeout} & \code{10} & Seconds without keyboard or mouse
      activity after which the values are saved and the window closes;
      \code{None} disables the automatic close. \\
    \bottomrule
  \end{tabularx}
\end{table}

\paragraph{Window behaviour.}
\begin{itemize}
  \item \emph{Save Parameters} validates all fields and writes a new file.
    Invalid values are listed in an error dialog and nothing is saved.
  \item \emph{OK} saves and closes the window; if the values are invalid the
    window stays open.
  \item \emph{Auto-save} (check box) saves every \code{save\_delay}
    milliseconds while enabled; problems are shown in the status line instead
    of dialogs.
  \item After \code{inactivity\_timeout} seconds without activity the window
    saves (if the values are valid) and closes. The default of ten seconds
    suits scripts that open the window at start-up and should continue
    unattended; pass a larger value or \code{None} for interactive editing.
  \item \emph{Browse} selects another output directory and \emph{Reset to
    Default} returns to the initial one. Closing the window with the window
    manager does not save.
\end{itemize}

\code{tune()} blocks until the window closes and returns the parser. The
values of the last \emph{successful} save (manual, OK, auto-save or
inactivity save) become the parser's defaults; if nothing was saved, the
parser is unchanged. The path of the last file is available as
\code{tuner.last\_saved\_path}. The class method
\code{ParameterTuner.tune\_parameters(parser, ...)} creates a tuner, calls
\code{tune()} and returns the parser; \code{ParameterAdjuster} is a
backward-compatible alias of the class. The font of the window is the class
attribute \code{font\_family} (default \code{"Times New Roman"}).

% norun
\begin{pycode}
from toolkit.parakit import ParameterTuner

tuner = ParameterTuner(parser, save_path="runs/params", inactivity_timeout=None,
                       validation_callbacks=validators)
parser = tuner.tune()                  # blocks until the window is closed
args = parser.parse_args()
print("saved to", tuner.last_saved_path)
\end{pycode}

\paragraph{Systems without Tk or without a display.}
\code{tune()} raises \code{ImportError} with installation instructions when
Tk is missing (the message is available as
\code{toolkit.parakit.TK\_MISSING\_MESSAGE}) and \code{RuntimeError} when no
display is available, for example on a cluster node or over SSH without X
forwarding. Scripts that should run everywhere catch both exceptions and fall
back to the headless API:

% norun
\begin{pycode}
try:
    ParameterTuner(parser, save_path="runs/params").tune()
except (ImportError, RuntimeError) as exc:
    print(f"Editor not available ({exc}); using the current defaults.")
\end{pycode}

The script \code{examples/parakit\_demo.py} demonstrates the complete
workflow, including the window (\code{--gui}), loading a file or directory
(\code{--load}) and command-line overrides (\code{--set KEY=VALUE}).
